\documentclass[10pt,a4paper]{article}

\usepackage[T2A]{fontenc}
\usepackage[utf8]{inputenc}
\usepackage[russian]{babel}
\usepackage{amsmath, amssymb, mathtools}

\usepackage{icomma}

\usepackage{geometry}
\geometry{a4paper, margin=10mm}

\begin{document}

    \section{Балльно-рейтинговая система оценки успеваемости}

    \subsection{Шкала перевода итогового балла в оценку}
    \begin{equation}
        \text{Итоговая оценка} =
        \begin{cases}
            \text{Неудовлетворительно}, & \text{при } B \in [0,\,54], \\
            \text{Удовлетворительно},   & \text{при } B \in [55,\,69], \\
            \text{Хорошо},              & \text{при } B \in [70,\,84], \\
            \text{Отлично},             & \text{при } B \in [85,\,100],
        \end{cases}
    \end{equation}
    где $B$ --- итоговый балл по дисциплине.

    \subsection{Расчет итогового балла}
    Итоговый балл определяется как взвешенная сумма компонентных показателей с последующим округлением в большую сторону:
    \begin{equation}
        B =
        \left\lceil
        0,3 \cdot \text{ЛР} +
        0,24 \cdot \text{КР} +
        0,1 \cdot \text{Посещаемость} +
        \text{Доп. баллы} +
        0,36 \cdot \text{Э}
        \right\rceil,
    \end{equation}
    где $\lceil \cdot \rceil$ --- функция потолка (округление до ближайшего целого в сторону увеличения),
    $ \text{ЛР} $ --- баллы за лабораторные (практические) работы,
    $ \text{КР} $ --- баллы за контрольные работы,
    $ \text{Посещаемость} $ --- баллы за посещение занятий,
    $ \text{Доп. баллы} $ --- дополнительные и компенсационные баллы за просроченные задания,
    $ \text{Э}$ --- баллы за экзамен.

    \subsection{Компонентные показатели}

    \subsubsection{Лабораторные (практические) работы (ЛР):}
    \begin{equation}
        \text{ЛР} = \frac{1}{N_{\text{ЛРобщ}}} \sum_{i=1}^{N_{\text{ЛРобщ}}} S_{\text{ЛР}}^{(i)},
    \end{equation}
    где $N_{\text{ЛРобщ}}$ --- общее количество лабораторных работ, $S_{\text{ЛР}}^{(i)}$ --- балл за $i$-ю работу.
    Невыполненные работы учитываются с $S_{\text{ЛР}}^{(i)} = 0$.

    Балл за отдельную лабораторную работу рассчитывается по формуле:
    \begin{equation}
        S_{\text{ЛР}}^{(i)} = 0,3 \cdot w_{\text{загр}}^{(i)} \cdot S_{\text{загр}}^{(i)} + 0,7 \cdot w_{\text{защ}}^{(i)} \cdot S_{\text{защ}}^{(i)},
    \end{equation}
    где $S_{\text{загр}}^{(i)}, S_{\text{защ}}^{(i)} \in [0;\,100]$ --- баллы за загрузку материалов в СДО и защиту работы соответственно.
    Коэффициенты своевременности определяются как:
    \begin{equation}
        w_{\text{загр}}^{(i)} = \min\!\left(1;\, \frac{W_{\text{план}}^{(i)}}{W_{\text{факт загр.}}^{(i)}}\right), \quad
        w_{\text{защ}}^{(i)} = \min\!\left(1;\, \frac{W_{\text{план}}^{(i)}}{W_{\text{факт защ.}}^{(i)}}\right),
    \end{equation}
    где $W_{\text{план}}^{(i)}$ --- номер недели планового дедлайна, $W_{\text{факт загр.}}$, $W_{\text{факт защ.}}$ --- номер недели фактического выполнения и защиты работы соответственно.
    Использование $\min(1;\,\cdot)$ исключает начисление баллов свыше $100$ при ранней сдаче.

    \subsubsection{Контрольные работы (КР):}
    \begin{equation}
        \text{КР} = \frac{1}{N_{\text{КРобщ}}} \sum_{j=1}^{N_{\text{КРобщ}}} S_{\text{КР}}^{(j)},
    \end{equation}
    где $N_{\text{КРобщ}}$ --- общее количество контрольных работ, $S_{\text{КР}}^{(j)} \in [0;\,100]$ --- балл за $j$-ю контрольную работу.

    \subsubsection{Посещаемость:}
    \begin{equation}
        \text{Посещаемость} = \frac{N_{\text{пос}}}{N_{\text{зан}}} \times 100,
    \end{equation}
    где $N_{\text{пос}}$ --- количество посещенных учебных занятий, $N_{\text{зан}}$ --- общее количество учебных занятий.

    \subsubsection{Экзамен (Э):}
    \begin{equation}
        \text{Э} = \sum_{k=1}^{N_{\text{зад}}} w_{\text{зад}}^{(k)} \cdot S_{\text{зад}}^{(k)},
    \end{equation}
    где $N_{\text{зад}}$ --- количество заданий (вопросов) экзамена, $w_{\text{зад}}^{(k)} \in (0;\,1]$ --- коэффициент за задание (вопрос), $\sum\limits_{k=1}^{N} w_{\text{зад}}^{(k)} = 1$, $S_{\text{зад}}^{(k)} \in [0;\,100]$ --- стоимость (фактический балл) выполнения $k$-го задания (ответа на вопрос).

    \subsection{Дополнительные баллы}

    \subsubsection{Варианты активностей для добора баллов}

    \begin{enumerate}
        \item \textbf{Расширенные лабораторные работы.}
        Студенту предлагается выполнить работу повышенной сложности или с расширенным объемом исследования.
        \begin{equation}
            S_{\text{расшЛР}}^{(i)} = \min\!\left(100;\, 1,2 \cdot S_{\text{баз}}^{(i)}\right),
        \end{equation}
        где $S_{\text{баз}}^{(i)} \in [0;\,100]$ --- балл за базовую версию работы, коэффициент $1,2$ отражает повышенную трудоемкость.

        \item \textbf{Аналитические эссе или обзоры литературы.}
        Письменная работа по теме пропущенного занятия.
        \begin{equation}
            S_{\text{эссе}} = \left(0,4 \cdot S_{\text{содерж}} + 0,3 \cdot S_{\text{структ}} + 0,3 \cdot S_{\text{оформл}}\right),
            \end{equation}
        где $S_{\text{содерж}}$, $S_{\text{структ}}$, $S_{\text{оформл}} \in [0;\,100]$ --- баллы за содержание, структуру и оформление работы соответственно.

        \item \textbf{Решение дополнительных задач.}
        Набор из 3--5 задач, превышающих базовый уровень сложности.
        \begin{equation}
            S_{\text{задачи}} = \frac{1}{N_{\text{задачи}}} \sum_{m=1}^{N} S_{\text{задачи}}^{(m)},
        \end{equation}
        где $ S_{\text{задачи}}^{(m)} \in [0;\,100]$ --- баллы за решение отдельной работы, $N_{\text{задачи}}$ --- количество задач в варианте.

        \item \textbf{Научно-исследовательская работа}
        Участие в научных мероприятиях, сборе данных, проведении экспериментов, оформлении отчетной документации по теме кафедры.
        \begin{equation}
            S_{\text{НИР}} = k_{\text{вклад}} \cdot S_{\text{НИРобщ}},
        \end{equation}
        где $S_{\text{НИРобщ}}$ --- максимальный балл, установленный для данной активности (как правило, $100$); $k_{\text{вклад}}  \in (0\,1]$ --- коэффициент личного вклада студента, определяемый научным руководителем в зависимости от объёма, сложности и качества проделанной работы.

        \item \textbf{Тестирование с открытым банком вопросов.} Сдача теста по материалам пропущенных лекций.
        \begin{equation}
            S_{\text{тест}} = \max\!\left(0;\, \frac{N_{\text{прав}} - 0,25 \cdot N_{\text{неправ}}}{N_{\text{общ}}} \times 100\right),
        \end{equation}
        где коэффициент $0,25$ --- штраф за угадывание.

        \item \textbf{Цифровые образовательные ресурсы (в том числе ДПО Политеха)}
        Прохождение онлайн-модулей, на платформе Мосполитеха или на внешней платформе (Stepik, GeekBrains и др.).
    \end{enumerate}

    \subsubsection{Ограничения и условия применения}

    \begin{itemize}
        \item Максимальная сумма баллов, полученных механизмами добора, не должна превышать $20\,\%$ от итоговой суммы компонент ЛР и КР:
        \begin{equation}
             \text{ДОП} = \sum S_{\text{доп}} \le 0,2 \cdot \left(0,3 \cdot \text{ЛР}_{\max} + 0,24 \cdot \text{КР}_{\max}\right).
        \end{equation}

        \item Не более двух механизмов добора могут быть применены к одной пропущенной теме.

        \item Общий срок выполнения компенсационных заданий не должен превышать 14 календарных дней с момента устранения причины пропуска.

        \item Инициатива применения механизма добора может исходить как от студента (возможно через письменное заявление), так и от преподавателя (рекомендация в индивидуальном плане).

        \item Все баллы, начисленные по механизмам добора, подлежат верификации: защита работы, устный опрос, проверка на оригинальность текста.

        \item Факт применения и результат добора фиксируются отдельно от основных баллов с пометкой <<компенсация>>.
    \end{itemize}

    \subsubsection{Исключительные ситуации}
    \begin{itemize}
        \item При пропуске по причине болезни сроком более 14 дней допускается индивидуальное согласование альтернативного графика добора баллов.

        \item Студентам, имеющим академические задолженности по другим дисциплинам, механизм добора может быть ограничен решением кафедры.

        \item Механизмы добора не применяются для компенсации пропусков итоговых аттестационных мероприятий (экзамен, зачет).
    \end{itemize}

    %%%%%%%%%%%%%%%%%%%%%%%%%%%%%%%%%%%%%%%%%%%%%%%%%%%%%%%%%%%%%%%%%%%%%%%%%%%%%

    \section{Вопросы для обсуждения}

    Для работоспособности системы необходимо по возможности как можно большее однообразие БРС в РПД. Поэтому необходимо обсудить на кафедре, и внести необходимые правки.

    \subsection{Баланс весовых коэффициентов между видами работ}

    \begin{enumerate}
        \item Насколько текущее распределение долей ($0,3$~ЛР, $0,24$~КР, $0,1$~Посещаемость, $0,36$~Экзамен) отражает фактическую трудоемкость и цели дисциплины?

        \item Целесообразно ли зафиксировать единую пропорцию для всей образовательной программы или предусмотреть дифференциацию весов в зависимости от профиля дисциплины (теоретический, прикладной, практико-ориентированный)?
    \end{enumerate}

    \subsection{Интеграция дополнительных форм учебной работы}
    \begin{enumerate}
        \item Следует ли расширить перечень оценочных средств за счет других видов работ (проектной деятельности, устных выступлений и т.д)?

        \item Каким образом обеспечить стандартизацию критериев оценки новых компонент, чтобы обеспечить единообразие на всех дисциплинах?

        \item Не приведет ли введение дополнительных компонент к общему их обесцениванию, снижению понимания формулы студентами и увеличению нагрузки на преподавателей?
    \end{enumerate}

    \subsection{Целесообразность дополнительного штрафа за выполнение работ}
    \begin{enumerate}
        \item В текущей редакции $\text{ЛР} = \frac{1}{N_{\text{ЛРобщ}}} \sum\limits_{i=1}^{N_{\text{ЛРобщ}}} S_{\text{ЛР}}^{(i)}$ при условии $S_{\text{ЛР}}^{(i)} = 0$ для невыполненных заданий уже реализует пропорциональное снижение. Требуется ли введение явного множителя $\frac{N_{\text{вып}}}{N_{\text{общ}}}$, приводящего фактически к двойному учету невыполненных работ?

        \item Мотивирует ли это студентов сдавать ВСЕ работы, или наоборот демотивирует обучающихся, выполнивших существенную часть заданий, но не сдавших отдельные работы?
    \end{enumerate}

    \subsection{Порядок учета пропусков и просрочек по уважительным причинам}
    \begin{enumerate}
        \item Какой алгоритм должен применяться при документально подтвержденных причинах: обнуление коэффициента своевременности ($k=1$), предоставление индивидуального графика сдачи или исключение пропущенного занятия из знаменателя формулы посещаемости?

        \item Каков регламентируемый срок предоставления подтверждающих документов, процедура их верификации преподавателем и порядок внесения корректировок в СДО?

        \item Следует ли установить количественный лимит на применение механизма уважительных причин за семестр для предотвращения злоупотреблений?

        \item Как применять штрафы за просрочку в случае РУП?
    \end{enumerate}

    \subsection{Порядок учета пропусков и просрочек по неуважительным причинам}

    \begin{enumerate}
        \item Следует ли предусмотреть механизм компенсации потерянных за счет пропуска/просрочки баллов по неуважительной причине? Если да, то в каком формате?
        Варианты:
        \begin{itemize}
            \item изначально предусмотренные в ЛР дополнительные задания;
            \item дополнительный банк заданий по дисциплине.
        \end{itemize}
        %Представляется логичным, ограничить количество дополнительных компенсационных баллов количеством потерянных
        \item Следует ли установить срок данной отработки, чтобы не столкнуться в конце семестра с огромным количеством студентов, желающих получить дополнительные задания? Или более эффективно будет сделать количество дополнительных заданий пропорциональным величин просрочки? Так будет не выгодно брать отработки за старые работы.
    \end{enumerate}

    \subsection{Документирование}
    \begin{enumerate}
        \item Готова ли инфраструктура СДО к автоматизированному расчету по данной формуле?
        \item Возможно заранее предусмотреть <<компенцационные>> задания для набора дополнительлных баллов.
        \item Возможно создать Excel документ с формулами. Тогда возникает проблема с <<хостингом>> для него.
    \end{enumerate}

    %%%%%%%%%%%%%%%%%%%%%%%%%%%%%%

\end{document}